\documentclass[10pt,letterpaper]{article}
\usepackage[letterpaper,margin=0.72in,headsep=0.22in]{geometry}
\usepackage[T1]{fontenc}
\usepackage[utf8]{inputenc}
\usepackage{lmodern}
\usepackage{xcolor}
\definecolor{accent}{HTML}{444444}
\usepackage{microtype}
\usepackage{enumitem}
\usepackage{titlesec}
\usepackage{hyperref}
\usepackage{xurl}
\usepackage{fancyhdr}
\usepackage{needspace}
\usepackage{setspace}
\usepackage{etoolbox}
\usepackage{fontawesome}
\setlength{\parindent}{0pt}
\setlength{\parskip}{0.22em}
\setlist[itemize]{itemsep=0.15em,topsep=0.15em,parsep=0pt,partopsep=0pt}
\urlstyle{same}
\hypersetup{colorlinks=true, urlcolor=black, linkcolor=black}
\emergencystretch=3em
\titleformat{\section}{\Large\sffamily\bfseries\color{accent}}{}{0em}{}[\vspace{0.15em}\color{accent}\titlerule]
\titlespacing*{\section}{0pt}{0.85em}{0.45em}
\titleformat{\subsection}{\normalsize\sffamily\bfseries}{}{0em}{}
\titlespacing*{\subsection}{0pt}{0.55em}{0.2em}
\renewcommand{\headrulewidth}{0.3pt}
\renewcommand{\footrulewidth}{0pt}
\setlength{\headheight}{13pt}
\usepackage[
  backend=biber,
  style=acmnumeric,
  sorting=none,
  defernumbers=true,
  maxbibnames=99,
  doi=true,
  url=true,
  isbn=false,
  eprint=false
]{biblatex}

\addbibresource{papers.bib}

% Count how many entries are in each refsection
\AtDataInput{%
  \csnumgdef{entrycount:\therefsection}{%
    \csuse{entrycount:\therefsection}+1}}
    
% Reverse the displayed bibliography numbers:
% newest printed item gets the highest number, oldest gets [1]
\DeclareFieldFormat{labelnumber}{\mkbibdesc{#1}}
\newrobustcmd*{\mkbibdesc}[1]{%
\number\numexpr\csuse{entrycount:\therefsection}+1-#1\relax}
\defbibenvironment{bibliography}
  {\list
     {\printtext[labelnumberwidth]{%
        \printfield{labelprefix}%
        \printfield{labelnumber}}}
     {\settowidth{\labelnumberwidth}{[\csuse{entrycount:\therefsection}]}%
      \setlength{\labelwidth}{\labelnumberwidth}%
      \setlength{\leftmargin}{\labelwidth}%
      \setlength{\labelsep}{\biblabelsep}%
      \addtolength{\leftmargin}{\labelsep}%
      \setlength{\itemsep}{\bibitemsep}%
      \setlength{\parsep}{\bibparsep}}%
   \renewcommand*{\makelabel}[1]{\hss##1}}
  {\endlist}
  {\item}
\newcommand{\bibliofont}{\small}
\setlength{\bibitemsep}{0.22\baselineskip}

\begin{document}


\begin{center}
{\LARGE\bfseries Kevin Buffardi, Ph.D.}\par
\vspace{0.25em}
{\small\href{mailto:kbuffardi@acm.org}{\faEnvelope~ kbuffardi@acm.org}
\enspace\textbar\enspace
\href{https://learnbyfailure.com}{\faGlobe~ learnbyfailure.com}
\enspace\textbar\enspace
\href{https://github.com/kbuffardi}{\faGithub~ github.com/kbuffardi}}
\end{center}

\needspace{2\baselineskip}\section*{Education} 
\textbf{Ph.D. in Computer Science}, Virginia Tech \\
\textbf{M.S. in Human-Computer Interaction}, with Distinction. DePaul University \\
\textbf{B.S. in Computer Science}, Cum Laude with Honors. University of Mary Washington

\section*{Experience}
\textbf{Full Professor}, 2024---Present; Associate Professor, 2019-2024; Assistant Professor, 2014-2019; \textit{Chico State}
\begin{itemize}[leftmargin=1.5em,itemsep=2pt, topsep=0pt]
\item Outstanding Professor Award (2022-2023), Chico State
\item Outstanding Teaching Award (2024) \textit{Nominated.}, Wang Excellence Award, CSU-wide excellence in teaching. 
\item Taught 17 courses---emphasis on Software Engineering, Usability, and K-12 Computer Science Education
\end{itemize}

\vspace{2mm}
\textbf{Dean's Teaching Fellow}, Computer Science Department; \textit{Virginia Tech}
\begin{itemize}[leftmargin=1.5em,itemsep=2pt, topsep=0pt]
\item Global Perspectives Scholar, Global Perspectives Program---Switzerland, 2013
\item Graduate Teaching Excellence Award, Virginia Tech, 2013
\item Outstanding Graduate Teaching Assistant, Computer Science Dept., 2012
\end{itemize}

\vspace{2mm}
\textbf{User Experience Specialist} 2006---2008; User Experience Analyst, 2005---2006; \textit{User Centric, LLC}, Chicago

\vspace{2mm}
\textbf{Usability Lab Intern} 2004; \textit{US Census Bureau}; Suitland, Maryland

\needspace{2\baselineskip}\section*{Research}
\subsection*{External Grants}
\begin{itemize}[leftmargin=1.5em,itemsep=2pt, topsep=0pt]
\item \textit{Collaborative Research: FRONTIERS: Fostering Software Engineering Competence through Valid, Reliable, and Practical Assessments of Individual Contributions to Team Projects}; NSF IUSE grant (\href{https://www.nsf.gov/awardsearch/show-award/?AWD_ID=2337271}{\#2337271} \href{https://www.nsf.gov/awardsearch/show-award/?AWD_ID=2337270}{\#2337270} \href{https://www.nsf.gov/awardsearch/show-award/?AWD_ID=2337269}{\#2337269}, \textbf{Awarded: \$750,000}, Co-PI, 2024---Present) 
\item \textit{Python 4 All}; California Education Learning Lab (\textbf{Awarded: \$50,000}, PI, 2024---2025) 
\item \textit{BPC-DP: Embedding Representation, Relevance, and Equity in a Personalized System of Instruction}; National Science Foundation, Broadening Participation in Computing (\href{https://www.nsf.gov/awardsearch/show-award/?AWD_ID=2315883}{\#2315883} \textbf{Awarded: \$299,912}, PI, 2023---2026) 
\item \textit{Coding Community: Inclusive Space for Programming Tutorials and Adaptive Learning}; California Education Learning Lab (\textbf{Awarded: \$100,000}, PI, 2020---2022) 
\item \textit{Coding Online Videos with Inclusive Demonstrations}; Department of Education American Rescue Plan Act Higher Education Emergency Relief Funds (\textbf{Awarded: \$36,834}, PI, 2021---2022) 
\item \textit{STEM Course Transformation: Cultivating a Culture of Entrepreneurial Mindset and Undergraduate Research}; NSF HSI (\href{https://www.nsf.gov/awardsearch/show-award/?AWD_ID=1953751}{\#1953751}, \textbf{Awarded: \$2,200,001}, Senior Personnel, 2020---2025) 
\end{itemize}

\subsection*{Peer-Reviewed Publications}
Published \textbf{45 papers}, concentrating on Software Engineering and Computer Science Education research \\
\textbf{14 mentored student researcher co-authors}, including \textbf{6 undergraduate co-authorships} \\
\textbf{\textit{Selected publications} from last 5 years}:
\begin{refsection}
\nocite{
hundhousen2026,
srivastava2026,
buffardi2026studentreflections,
more2026pareto,
bijoor2025gitdash,
wang2025expanding,
buffardi2025consistency,
buffardi2024designing,
buffardi2023cognitive,
buffardi2023codevid,
buffardi2022integrating,
buffardi2022codewitus,
morrison2022evidence,
buffardi2021relevant,
alavi2021redesigning,
buffardi2021unittestsmells,
}
\printbibliography[heading=none]
\end{refsection}


\needspace{2\baselineskip}\section*{Scholarship}
\begin{itemize}[leftmargin=1.25em]
\item Program Committee: International Conference on Software Engineering (ICSE-SEET), International Conference on Software Engineering Education and Training (CSEE\&T), International Conference on the Foundations of Software Engineering (FSE)
\item Invited NSF panelist: proposal reviewer for CISE directorate, 2017, 2022
\item Associate Editor: PLOS One (2021---Present);
\item Associate Editor: Frontiers in ICT, section on Digital Education (2021---2022)
\item Associate Program Chair: ACM Technical Symposium on Computer Science Education (SIGCSE), ACM Innovation and Technology in Computing Education (ITiCSE)
\item Reviewer ACM International Computer Education Research (ICER), ACM Transactions of Computing Education (TOCE), Journal on Empirical Software Engineering (EMSE), ACM DATA BASE for Advances in Information Systems (SIGMISDB), Computer Science Education, Taylor and Francis
\end{itemize}
\needspace{2\baselineskip}\section*{Relevant Professional Development \& Service}
\begin{itemize}[leftmargin=1.5em]
    \item AI curriculum committee group, \textit{Chair}, CSCI department, 2026
    \item Research and Sponsored Programs Grant Writing Bootcamp, Summer, 2018
    \item Hispanic Serving Institution (HSI) Faculty Learning Community, Spring, 2017
    \item Professor's Open Source Software Experience (POSSE) Foss2Serve (\url{http://foss2serve.org/}) Professional development for instructors interested in student participation in Humanitarian Free and Open Source Software
    \item Funded (awarded \$4000) development of open source course description and series of educational materials for a module on software testing and quality assurance
    \item Collaborated and spoke as part of a peer-reviewed, five-person panel (SIGCSE 2017) on community engagement in open source projects
    \item NSF Grant Writing Workshop, attended in Portland, Oregon, 2016
\end{itemize}

\end{document}
